本章将介绍 C++11 中与继承（inheritance）以及从其他构造函数调用构造函数相关的改进。

\section{委托构造函数}

有时，当你的类包含许多数据成员和多个构造函数时，复用它们的初始化代码会很方便。幸运的是，从 C++11 起，你可以使用委托构造函数（delegating constructor）。让我们看一个例子：

\begin{codeboxt}[cpp]{代码清单 4.1 委托构造函数（在 Compiler Explorer 上运行）}
class Product {
public:
Product(int id, unsigned quantity, const std::string& name)
: id_{id}, quantity_{quantity}, name_{name} {
verifyData();
}
Product(const std::string& name)
: Product{0, 0, name}
{ }
void verifyData() {
if (quantity_ > MaxQuantity)
throw std::invalid_argument("quantity is too large!");
}
const std::string& getName() const { return name_; }
private:
int id_;
unsigned quantity_;
std::string name_;
static constexpr unsigned MaxQuantity = 100;
};
\end{codeboxt}

在上面的例子中，我们声明了两个构造函数。第一个负责核心工作，第二个调用那个“主要的”构造函数。在这个主构造函数中，我们不仅初始化数据成员，还调用了其他代码。在我们的例子中，这是一种基础的数据校验形式。请注意，我还为构造函数使用了一个默认参数（\texttt{id = 0}），所以当你想提供不同的调用选项时，这是另一种替代方案。

下面是演示代码：

\begin{codeboxt}[cpp]{代码清单 4.1 委托构造函数（演示）（在 Compiler Explorer 上运行）}
int main() {
try {
Product box{"a box"};
std::cout << "product: " << box.getName() << " created... \n";
Product toy{101, 200, "a box"};
std::cout << "product: " << toy.getName() << " created... \n";
}
catch (const std::exception& e) {
std::cout << "cannot create: " << e.what() << '\n';
}
}
\end{codeboxt}

运行后我们会得到以下输出：

\begin{codebox}[cpp]
product: a box created...
cannot create: quantity is too large!
\end{codebox}

如果没有委托构造函数，我们就得重复代码：

\begin{codebox}[cpp]
Product(int id, unsigned quantity, const std::string& name)
: id_{id}, quantity_{quantity}, name_{name} {
verifyData();
}
Product(const std::string& name, int id = 0)
: id_{id}, quantity_{0}, name_{name} {
verifyData();
// 代码重复
}
\end{codebox}

如你所见，使用委托构造函数的代码更加紧凑，并且可以实现完整的代码复用。这能节省打字量，并可能消除代码中的各种“copy\&paste”式错误。

此外，这种语法并不局限于普通的构造函数，你还可以从拷贝构造函数或移动构造函数中调用另一个构造函数：

\begin{codebox}[cpp]
// 拷贝：
Product(const Product& other) : Product{other.id_, other.quantity_, other.name_}
{ }
// 移动，可能是个坏主意（仅作示意）
Product(Product&& other)
: Product{other.id_, other.quantity_, other.name_}
{ }
\end{codebox}

对于拷贝构造函数，这样的代码可以复用校验部分。但要当心移动构造函数，因为上面的代码并不会产生任何“移动”，而是会拷贝数据，这就违背了它的主要目的。

\begin{myWarnBox}
要当心这种语法！

\begin{codebox}[cpp]
explicit PropertyInfo(double price) { PropertyInfo(...); }
\end{codebox}

上面这行会创建一个局部对象，而不是调用另一个构造函数！对构造函数的调用必须出现在构造函数体之前。
\end{myWarnBox}

\section{限制}

编写过多的构造函数可能会导致一些错误和递归调用。请看下面的代码：


\begin{codebox}[cpp]
class Product {
public:
Product(int id, unsigned quantity, const std::string& name)
: Product {name, id} { }
Product(const std::string& name, int id = 0)
: Product{id, 0, name} { }
// ...
};
Product recursion{"a single recursion"};
\end{codebox}

当 \texttt{info} 对象调用它的构造函数时会发生什么？

你可能会遇到段错误（segmentation fault）和栈溢出！这是一个递归调用，编译器在代码运行时之前无法检测到它。

另一个“限制”是，你不能把成员初始化与调用其他构造函数混在一起。

下面的代码无法编译：

\begin{codebox}[cpp]
Product(int id, unsigned quantity, const std::string& name)
: Product {name, id}, quantity_{quantity} { }
\end{codebox}

例如，GCC 报告以下错误：

\myGraphic{0.9}{images/fig_p082_1.png}{}

\begin{codebox}[cpp]
error: mem-initializer for 'Product::quantity_' follows constructor delegation
7 |
: Product {name, id}, quantity_{quantity}
\end{codebox}

总之，如果你想使用委托构造函数，就不能初始化其他数据成员。

接下来我们进入关于构造函数的另一节，你将学到 Modern C++ 的另一个技巧。

\section{继承}

我们来看看你的类继承自其他类的情况。构造函数会发生什么？编译器何时调用它们？这一讨论将为 C++11 中一项名为“继承构造函数（inheriting constructors）”的新特性提供背景。

为了调试，我们从 \texttt{DataPacket} 引入一个名为 \texttt{DebugData} 的派生类，它具有特殊的打印能力：

\begin{codebox}[cpp]
class DebugDataPacket : public DataPacket {
public:
DebugDataPacket(const std::string& name, size_t serverId)
: DataPacket{name, serverId} {
}
void DebugPrint(std::ostream& os) {
os << getData() << ", " << getCheckSum() << '\n';
}
};
\end{codebox}

如你所见，代码声明了一个新类，并使用 \texttt{: public DataPacket} 来表示公有继承。该例子还定义了一个调用基类构造函数的单一构造函数。

C++ 提供了三种方式来指定我们从基类继承的方式：

• 公有（public）继承意味着基类的 \texttt{public} 成员成为派生类的 \texttt{public} 成员，而 \texttt{protected} 成员在派生类型中仍是 \texttt{protected}。

• 保护（protected）继承使基类的所有 \texttt{public} 和 \texttt{protected} 成员都能作为派生类的 \texttt{protected} 成员访问。

• 私有（private）继承使基类的所有 \texttt{public} 和 \texttt{protected} 成员都能作为派生类的 \texttt{private} 成员访问。

• 在这三种情况下，基类的 \texttt{private} 成员对派生类不可访问，除非显式地声明为 \texttt{friend}。

我们可以这样使用它：

\begin{codeboxt}[cpp]{代码清单 4.2 继承，简单演示代码（在 Compiler Explorer 上运行）}
int main() {
DebugDataPacket hello{"hello!", 404};
hello.DebugPrint(std::cout);
}
\end{codeboxt}

在这个例子中，基类构造函数是显式调用的，但一般而言，每个构造函数还会隐式地调用基类的默认构造函数。下面的代码说明了这一点：

\begin{codeboxt}[cpp]{代码清单 4.3 基类构造顺序（在 Compiler Explorer 上运行）}
#include <iostream>
#include <string>
class Product {
public:
Product() : id_{0} { std::cout << "Product() default\n"; }
explicit Product(int id, const std::string& name)
: id_{id}, name_{name} {
std::cout << "Product(): " << id_ << ", " << name_ << '\n';
}
protected:
int id_;
std::string name_;
};
class ExProduct : public Product {
public:
ExProduct() { std::cout << "ExProduct() default\n"; }
explicit ExProduct(int id) {
id_ = id;
std::cout << "ExProduct(id)\n";
}
};
int main() {
ExProduct p;
ExProduct withId{42};
}
\end{codeboxt}

如果运行程序，我们会得到以下输出：

\begin{codebox}[cpp]
Product() default
ExProduct() default
Product() default
ExProduct(id)
\end{codebox}

如你所见，即使我们没有在 \texttt{ExProduct} 构造函数中调用任何基类构造函数，编译器仍然调用了它。更重要的是，在派生类的构造函数中，你不能在初始化列表中直接使用基类的数据成员，例如：

\begin{codebox}[cpp]
ExProduct(): id_(10) { // << 错误！我们没有访问权限！
std::cout << "ExProduct() default\n";
}
\end{codebox}

你只能在构造函数体中访问它：

\begin{codebox}[cpp]
ExProduct(int id) {
id_ = id;
std::cout << "ExProduct(id)\n";
}
\end{codebox}

这种行为对保持对象的完整性至关重要。

\begin{myWarnBox}
此外，最好不要在构造函数中调用 \texttt{virtual} 函数，因为它们的行为可能与预期不同。简而言之，在基类构造函数中调用虚函数会调用基类的实现，因为派生类及其实现尚未建立。你可以在 C++ FAQ\footnote{\url{https://isocpp.org/wiki/faq/strange-inheritance\#calling-virtuals-from-ctors}} 或 C++ Core 指南 C.82\footnote{\url{https://isocpp.github.io/CppCoreGuidelines/CppCoreGuidelines\#Rc-ctor-virtual}} 中阅读更多关于此行为的内容。
\end{myWarnBox}


在介绍完继承主题之后，我们可以讨论 Modern C++ 带来的一个改进。

\section{继承构造函数}

在先前的 \texttt{DebugPropertyInfo} 例子中，我们没有任何新的数据成员，只有一些新的成员函数。代码中展示了一个调用基类构造函数的单一构造函数。从 C++11 起，你可以告诉编译器“复用”这段代码：

\begin{codeboxt}[cpp]{代码清单 4.4 继承构造函数（在 Compiler Explorer 上运行）}
class DebugDataPacket : public DataPacket {
public:
using DataPacket::DataPacket;
void DebugPrint(std::ostream& os) {
os << getData() << ", " << getCheckSum() << '\n';
}
};
int main() {
DebugDataPacket hello{"hello!", 404};
hello.DebugPrint(std::cout);
}
\end{codeboxt}

\subsection{考虑第 3 行：}

\begin{codebox}[cpp]
using DataPacket::DataPacket;
\end{codebox}

这告诉编译器它可以使用基类的所有构造函数，忽略访问修饰符。这意味着所有公有构造函数都可见且可被调用，但受保护的构造函数在那个上下文中仍是 \texttt{protected}。看下面的例子：

\begin{codeboxt}[cpp]{代码清单 4.5 继承构造函数与保护段（在 Compiler Explorer 上运行）}
struct Base {
int x{};
int y{};
Base(int a, int b): x{a}, y{b} { }
protected:
Base() = default;
Base(int a): x{a} { }
};
struct Derived : public Base {
using Base::Base;
};
int main() {
// Derived d{0};
// 错误：'Base::Base(int)' 是受保护的
Derived d2{0, 1}; // 正确
}
\end{codeboxt}

如果你想限制对构造函数的访问，就必须为 \texttt{Derived} 显式地编写构造函数：

\begin{codebox}[cpp]
Derived(int a) : Base{a} { }
Derived d{0}; // 现在正确，因为 Derived::Derived(int) 是公有的
\end{codebox}

我们讲完了关于构造函数的所有内容，但最好再提一件事：析构函数。详见下一章。
